\section{Materiality 与 Disclosure：四工作日从哪里开始算}

安全团队可以判断技术影响，却不能独自决定全部披露义务。\term{materiality}（重大性）在美国证券法语境中关注一个合理投资者是否很可能认为该信息对投资决定重要；它不是由记录数量、CVSS、停机分钟或攻击者标签单独决定。不同司法辖区、行业、合同和用户通知规则还有各自 trigger，因此组织需要 parallel analysis，而不是寻找一个“万能阈值”。

\subsection{Facts、materiality、audience 与 timing}

披露工作流先建立可靠事实，再由适当人员分析对运营、财务、声誉、法律和战略的影响，然后识别受众与时间。受众可能包括受影响用户、业务伙伴、保险人、执法机构、隐私或行业监管者，以及公开市场。不同通知可以有不同目的和信息粒度，但不能互相矛盾。读这一流程时要特别注意：technical severity 只回答防御资源应多快投入，materiality 则回答该事件对特定法律和决策受众是否重要；前者可以触发后者，却不能自动替代后者。

\lecturefigure{08-disclosure-decision.png}{披露决策从技术事实开始，经 materiality 分析进入受众、内容与时间选择。}{SEC 2023 cybersecurity disclosure final rule；NIST guidance；概念重绘。}

\begin{knowledgebox}{读图：先判断，再启动相应时钟}
Facts 列出 scope、data、systems、operational impact 与 confidence；materiality 分析 investor/customer/user harm、financial/operational effect 和 cumulative impact；audience 列出 users、regulators、law enforcement、partners；timing 再映射 contract、statute、Form 8-K 或 approved delay。关键顺序是“发现后不无理拖延地作重大性判断；一旦确定重大，再计算相应申报窗口”，不是所有事件发现后统一四天。
\end{knowledgebox}

\begin{importantbox}{SEC 四工作日规则的准确口径}
对适用的美国上市公司，SEC 规则要求在发现 incident 后不无理拖延地判断重大性；如果公司确定事件具有重大性，通常在该\textbf{重大性确定之后四个工作日内}提交 Form 8-K Item 1.05。\term{Form 8-K} 是上市公司报告特定重大事件的当前报告表格。规则包含由美国司法部长认定即时披露会对国家安全或公共安全造成重大风险时的延迟机制。具体适用必须由合格 counsel 判断。
\end{importantbox}

\begin{warningbox}{不要等待“所有事实百分之百确定”}
调查可能持续数周，但通知时钟可能更短。治理目标不是用不确定性拖延，而是明确已知、未知、置信度和预计更新。过早给出不准确数字会伤害信任，过晚启动 materiality workstream 也会制造违规风险。应从事件早期就让 disclosure owner 参与，并设置重复评估 checkpoint。
\end{warningbox}

\begin{knowledgebox}{RACI 如何防止责任扩散}
\term{RACI} 是 Responsible、Accountable、Consulted、Informed 的责任矩阵。对 technical scoping，forensics lead 可能 Responsible、IC Accountable；对法律义务，counsel Responsible、GC Accountable；对重大经营披露，CEO 或 disclosure committee Accountable。RACI 不是让四类角色平均承担责任，而是明确每个 deliverable 的唯一最终责任人。
\end{knowledgebox}

\subsection{本章小结}

Disclosure 是证据驱动、跨职能且多时钟的决策。安全团队提供及时可靠事实，counsel 解释义务，高管或委员会承担经营与公开披露决定；“四天”从重大性确定开始，而不是从第一个告警开始机械倒计时。

